\documentclass[a4paper,11pt,twocolumn]{article}
\usepackage{graphicx} % Required for inserting images
\usepackage[utf8]{inputenc} % poza overleaf wymagane
\usepackage{polski} % polskie znaki
\usepackage{hyperref}
\usepackage{booktabs}
\usepackage{ragged2e} % for \justify
\usepackage{listings}

\title{Notatka Latex}
\author{Moje imię i nazwisko}
\date{12 Marca 2026}

\begin{document}

\twocolumn[
  \maketitle
  \vspace*{-5em}
  \section*{}
  \centering
  
    \begin{minipage}{0.88\textwidth}  % Adjust the width as needed (e.g., 0.8\textwidth)
      \justify
      Przykładowe streszczenie poza trybem dwukolumnowym.
      The study of magnetic domains in data carriers was conducted, leading to a deeper understanding of Atomic Force Microscopy (AFM) principles and greater familiarity with the Magnetic Force Microscopy (MFM) method. 
    \end{minipage}
  
  \vspace*{2em}
]

\section{Podstawowe elementy}
\begin{itemize}
        % \centering
        \item \textbf{section\{...\}}, \textbf{subsection\{...\}},  \textbf{subsubsection\{...\}}, \textbf{*section\{...\}} - tworzą odpowiednio sekcje, podsekcje oraz sekcje nienumerowaną.
        \item \textbf{itemize}, \textbf{enumerate} - tworzą listy punktowane i~numerowane. Rozpoczyna się je od \textbf{begin} i kończy \textbf{end}. Każdy element musi być poprzedzony poleceniem \textbf{item}.
        \item \textbf{flushleft}, \textbf{flushright}, \textbf{center} lub \textbf{centering} - wyrównują tekst odpowiednio do lewej, prawej lub do środka. Różnica między nimi polega na tym, że pierwsze trzy są środowiskami (rozpoczynają się od \textbf{begin} i kończą \textbf{end}), podczas gdy \textbf{centering} jest poleceniem, które można użyć w dowolnym miejscu w dokumencie.
\end{itemize}

\section{Równania}
\subsection{Przykładowe}
\begin{equation}
    w(z) = w_0 \sqrt{1+ M^2(\frac{z-z_0}{z_R})^2}
\end{equation}

\section{Tabele}
\subsection{Przykładowa}


\begin{table}[htp]
  \centering
  \caption{Podpis na górze}
  \label{tab:nazwaTabeli}
  \scalebox{1.0}{
  \begin{tabular}{ccc}
    \toprule
    \textbf{$z_0$ [mm]} & \textbf{$w_0$ [mm]} & \textbf{$M^2$ [-]} \\
    \midrule
    $552.627 \pm 0.047$ & $0.0115 \pm 0.0086$ & $ 2.0 \pm 3.0$  \\
    \bottomrule
  \end{tabular}}
\end{table}



\section{Obrazy}
\subsection{Przykładowa grafika}
\begin{figure}[htp]
    \centering
    \includegraphics[width=0.4\textwidth]{stsci-01gnymeg4ym0ksh1x63g894t6h.png}
    \caption{Podpis pod obrazem}
    \label{fig:nazwaObrazu}
\end{figure}

\section{Lista}
\subsection{Przykładowa lista Linux}
\begin{lstlisting}[language=bash, caption={Przykładowa lista poleceń Linux}, label={lst:linux}]

#!/bin/bash

# 1. communication
echo "this is a script"


\end{lstlisting}

\begin{thebibliography}{}

    \bibitem{Klucz}
    \href{}{Nazwisko Imię, \textit{Tytuł italic}, wydawnictwo, rok.}

\end{thebibliography}

\end{document}
